% Results update: E1/E2 stage localisation and robustness
% Paste into Chapter 4 after the main E1/E2 result paragraphs.

\subsection{Robustness of stage localisation}
\label{sec:e1-stage-robustness}

The initial E1 analysis identified a reproducible early reorganisation of weight space across Pythia model sizes. I next asked whether this window was an artefact of a single plot, a single metric family, or a small number of layers. To make the phase localisation more explicit, I re-analysed the E1 metrics as an adjacent-checkpoint transition problem. For each model size, module, layer, and spectral indicator, I computed the adjacent interval with the largest transition strength, and then aggregated these interval-level votes across the model.

The resulting localisation is consistent across the larger Pythia models. For Pythia-160M, 410M, and 1B, the strongest aggregate boundary vote is the interval from 1000 to 2000 steps. The neighbouring 512--1000 interval is the second strongest interval across the same models, so the most defensible interpretation is not a single instantaneous boundary but a short early developmental window concentrated between approximately 512 and 2000 steps. This window is also consistent with the earlier E1 visual analysis: stable-rank collapse, dominant spectral-direction formation, and leading-subspace instability all concentrate in the first few thousand steps, long before the final checkpoint.

Layer-bootstrap analysis strengthens this conclusion. When layers are resampled with replacement and the boundary vote is recomputed, the 1000--2000 interval remains the top interval in nearly all bootstrap samples for Pythia-160M, Pythia-410M, and Pythia-1B. This means that the localisation is not explained by a single unusual layer. Metric-family ablations give a similar conclusion: the 1000--2000 interval remains dominant when rank metrics, MP-outlier metrics, or subspace-stability metrics are removed. The one important qualification is that different metric families appear to detect different parts of the transition. MP-outlier and subspace-stability metrics tend to fire earlier, around 128--512, while rank-compression metrics and the multi-metric aggregate peak around 1000--2000. I therefore treat the E1 result as a staged reorganisation rather than a point transition: early subspace destabilisation and outlier emergence are followed by broader rank compression and consolidation.

The E2 functional-grounding results show broad temporal overlap with this E1 window, but not a tight interval-by-interval correspondence. Fixed-text next-token loss improves rapidly in the earliest checkpoints, while syntax-regularity margins rise mainly across the 128--2000 step interval. Arithmetic and sequence-completion probes remain noisy and exploratory. Thus, E2 supports the claim that the E1 window is behaviourally consequential at a coarse level, but the phase boundary itself remains defined by the independent weight-space analysis rather than by behavioural probes.

Overall, E1 and its robustness analyses justify the intervention grid used in E3. The grid samples before the window, inside the window, immediately after the aggregate boundary, and in late consolidated checkpoints. Importantly, the intervention experiment does not fit the phase boundary from the durability outcome. The candidate window is fixed by E1, and E3 asks whether independently measured durability changes across that window.
